% 带权反例：边 {s,t} 长5，{s,a}、{a,t} 各长1。BFS 以1条边到达t（成本5），s-a-t 用2条边、成本2。
\begin{tikzpicture}[font=\figurefont\fontsize{8}{10}\selectfont,
  vtx/.style={circle,draw=FigureStroke,line width=.8pt,minimum size=6.5mm,inner sep=0pt,text=FigureInk,fill=white},
  wt/.style={fill=white,inner sep=1pt,text=FigureInk},
  note/.style={font=\FigureNoteFont\fontsize{7.5}{9}\selectfont,text=FigureMuted}]
  \node[vtx,fill=FigureYellowFill] (s) at (0,0) {$s$};
  \node[vtx] (a) at (1.6,1.1) {$a$};
  \node[vtx] (t) at (3.2,0) {$t$};
  \draw[FigureBlue,line width=1.3pt] (s)--node[wt,below]{$5$}(t);
  \draw[FigureGreen,line width=1.3pt,dashed] (s)--node[wt,above left]{$1$}(a);
  \draw[FigureGreen,line width=1.3pt,dashed] (a)--node[wt,above right]{$1$}(t);
  \node[anchor=west,text=FigureBlue] at (4.0,0.35) {BFS 路线 $s\text{--}t$：1 条边，总长 5};
  \node[anchor=west,text=FigureGreen] at (4.0,-0.25) {另一路线 $s\text{--}a\text{--}t$：2 条边，总长 2};
  \node[anchor=west,note] at (4.0,-0.85) {边数最少不等于总长最小};
\end{tikzpicture}
